% page 382 of the facsimile (image p-381 of the scan)
% block 157.5 x 216.0 mm ; left margin 8.0 mm
\pgc{-12.700mm}{0.000mm}{
\l{\cel{3}374}
\l{}
\l{\sou{min}.\cel{2}Adverbo qua indikas la inferioreso pri qualeso, quan-}
\l{\cel{3}teso, precozeso, extenseso, e c.\cel{2}= ne tam multe ....-a.}
\l{\cel{2}- FIRS.}
\l{}
\l{minar. (trans.) Exkavar truo en qua onu pozas pulvero quan}
\l{\cel{3}onu acendas por destruktar exploze to quo esas sur o super}
\l{\cel{3}lu. - DEFIRS.}
\l{}
\l{\sou{min}acar. (\sou{trans., per}). I. Manifestar ad (ulu) sua intenco}
\l{\cel{3}nocar lu. Avertar (ulu) ke lu timez ulo de altru. - II.}
\l{\cel{3}(\sou{aludante kozo}).Manifestar, per kelka indici, ke, de lu,}
\l{\cel{3}esas timenda ulo kom eventonta balde. - EFIS.}
\l{}
\l{\sou{minareto.}\cel{2}Turmo di moskeo. - DEFIRS.}
\l{}
\l{\sou{mineralo}.\cel{2}Substanco qua konsistas ye materio kruda, neorga-}
\l{\cel{3}nizita, neorganika. DEFIRS.}
\l{}
\l{\sou{mineralogio}. Ta parto di la natur-historio qua traktas la}
\l{\cel{3}minerali, la korpi ne-organika tala qualan kande on renkon-}
\l{\cel{3}li en tero o sur la surfaco di olca. - DEFIRS.\cel{6}tras}
\l{}
\l{\sou{minestrelo}.(en\cel{1}\sou{la mezepoko}).\cel{2}Poeto, muzikisto recitanta,}
\l{\cel{3}qua kantis sua kompozuri ed olti di la poeti cetera.}
\l{\cel{3}- DEFIR.}
\l{}
\l{\sou{miniaturo}. I. Litero reda, trasita per miniumo, sur la ma-}
\l{\cel{3}nuskripti, la meso-libri, e c. por ornar la komenco-loko}
\l{\cel{3}di la chapitri. - Litero ye kolori diversa. - II. Pikturi}
\l{\cel{3}delikata, ek mikra temi, sur velino, pergamo, en la ma-}
\l{\cel{3}nuskripti, la meso-libri. - III. Pikturo tre delikata,}
\l{\cel{3}tre mikra, facita per farbi dilutita en aquo e gumo,}
\l{\cel{3}sur ivoro, sur velino. - IV. Reprezenturo tre mikra-dimen-}
\l{\cel{3}siona. - DEFIRS.}
\l{}
\l{\sou{minim}. (\sou{gram.)}\cel{2}Superlativo absoluta di \textquotedbl{}pokeso\textquotedbl{}. - DEFIRS.}
\l{}
\l{\sou{miniona}.Di qua la mikreso esas charmoza. F.}
\l{}
\l{\sou{ministerio}. Laboreyo di ministro e di la oficieri ed komizi}
\l{\cel{3}di lua departmenti. (EF.)}
\l{}
\l{\sou{ministro}. Viro (o muliero) a qua la povo leg-aplikanta kon-}
\l{\cel{3}fidis un ek la taski precipua di la guverno. DEFIRS.}
\l{}
\l{\sou{miniumo}. I. (\sou{kemio)}. Deutoxido di plumbo, substanco farba,}
\l{\cel{3}reda. - II. Farbo oleoza, facita ek miniumo, uzata kom}
\l{\cel{3}strato-unesma lor la induto di fero. - DEFIRS.}
}
